% !TeX root = part2.tex

\documentclass[../final.tex]{subfiles}

\begin{document}

% При отдельной сборке продолжаем нумерацию полного конспекта.
\ifSubfilesClassLoaded{\setcounter{chapter}{1}\setcounter{section}{5}}{}


% ============================================================
\section{Потери энергии электронов}
% ============================================================

Для электронов необходимо учитывать, что налетающая частица
и электрон вещества являются одинаковыми частицами.
В результате энергетический спектр выбитых электронов отличается
от случая тяжёлых заряженных частиц.

Дифференциальное число взаимодействий имеет вид

\begin{equation*}
    \frac{d^2n}{d\varepsilon\,dx}
    =
    C
    \frac{E^2}{(E-\varepsilon)^2\varepsilon^2}
    \left[
        1-\frac{\varepsilon}{E}
        +\left(\frac{\varepsilon}{E}\right)^2
    \right].
\end{equation*}

Здесь $E$ --- энергия налетающего электрона,
а $\varepsilon$ --- энергия, переданная при взаимодействии.

Для описания энергетической зависимости потерь удобно использовать
лоренц-фактор

\begin{equation*}
    \gamma=\frac{E}{m}.
\end{equation*}

При этом для тяжёлых заряженных частиц, электронов и позитронов
характер энергетических потерь различается.

\begin{rbox}[left=18pt,right=18pt,top=14pt,bottom=10pt]{[]}
    \centering

    \begin{tikzpicture}[
        x=1cm,
        y=1cm,
        >=Latex,
        every node/.style={text=rtext}
    ]

        % ------------------------------------------------------------
        % Оси
        % ------------------------------------------------------------

        \draw[
            ->,
            draw=rtext,
            line width=0.9pt
        ]
        (0,0) -- (9.0,0);

        \draw[
            ->,
            draw=rtext,
            line width=0.9pt
        ]
        (0,0) -- (0,5.0);


        % ------------------------------------------------------------
        % Деления по оси X
        % Логарифмический масштаб:
        % 10, 100, 1000, 10000 расположены равномерно
        % ------------------------------------------------------------

        \foreach \x/\lab in {
            1.25/10,
            3.15/100,
            5.05/1000,
            6.95/10000
        }{
            \draw[
                draw=rtext,
                line width=0.75pt
            ]
            (\x,-0.10) -- (\x,0.10);

            \node[
                below=4pt,
                font=\small
            ]
            at (\x,0)
            {$\lab$};
        }


        % ------------------------------------------------------------
        % Деления по оси Y
        % ------------------------------------------------------------

        \draw[
            draw=rtext,
            line width=0.75pt
        ]
        (-0.10,0.80) -- (0.10,0.80);

        \node[
            left=4pt,
            font=\small
        ]
        at (0,0.80)
        {$1.5$};


        \draw[
            draw=rtext,
            line width=0.75pt
        ]
        (-0.10,2.15) -- (0.10,2.15);

        \node[
            left=4pt,
            font=\small
        ]
        at (0,2.15)
        {$2$};


        % ------------------------------------------------------------
        % Подписи осей
        % ------------------------------------------------------------

        \node[
            anchor=south,
            font=\small
        ]
        at (-0.05,4.75)
        {$\displaystyle \frac{dE}{dx}$};

        \node[
            below=18pt,
            font=\small
        ]
        at (4.3,0)
        {$\displaystyle \gamma=\frac{E}{m}$};


        % ------------------------------------------------------------
        % Пион
        %
        % Быстрое падение,
        % широкий минимум,
        % затем медленный релятивистский рост
        % ------------------------------------------------------------

        \draw[
            draw=rc3,
            line width=1.55pt
        ]
        (0.48,4.55)
        .. controls
            (0.52,3.55)
            and
            (0.72,2.75)
        ..
        (1.22,2.48)
        .. controls
            (1.85,2.12)
            and
            (2.85,1.98)
        ..
        (4.55,1.92)
        .. controls
            (5.20,1.91)
            and
            (5.95,2.08)
        ..
        (7.05,2.42);

        \node[
            font=\Large
        ]
        at (3.75,2.55)
        {$\pi$};


        % ------------------------------------------------------------
        % Электрон e^-
        % ------------------------------------------------------------

        \draw[
            draw=rc4,
            line width=1.45pt
        ]
        (1.55,0.92)
        .. controls
            (2.35,1.16)
            and
            (3.40,1.43)
        ..
        (4.45,1.67)
        .. controls
            (5.45,1.90)
            and
            (6.35,2.13)
        ..
        (7.55,2.42);


        % ------------------------------------------------------------
        % Позитрон e^+
        % ------------------------------------------------------------

        \draw[
            draw=rc2,
            line width=1.45pt
        ]
        (1.55,0.70)
        .. controls
            (2.35,0.97)
            and
            (3.40,1.22)
        ..
        (4.45,1.47)
        .. controls
            (5.45,1.70)
            and
            (6.35,1.93)
        ..
        (7.55,2.20);


        % ------------------------------------------------------------
        % Подписи электронных кривых
        % ------------------------------------------------------------

        \node[
            anchor=west,
            font=\large
        ]
        at (7.65,2.42)
        {$e^-$};

        \node[
            anchor=west,
            font=\large
        ]
        at (7.65,2.18)
        {$e^+$};

    \end{tikzpicture}

    \captionsetup{
        font=footnotesize,
        labelfont=bf,
        skip=5pt
    }

    \captionof{figure}{
        Качественная зависимость удельных ионизационных потерь
        от лоренц-фактора
        $\gamma=E/m$
        для пиона, электрона и позитрона.
    }
    \label{fig:ionization_losses_particles}
\end{rbox}


% ------------------------------------------------------------
\subsection{Пробег частицы}
% ------------------------------------------------------------

Если известна зависимость энергетических потерь от энергии,
то полный пробег частицы можно определить интегрированием:

\begin{equation*}
    R
    =
    \int_0^{E_0}
    \frac{1}{dE/dx}\,dE,
\end{equation*}

где $E_0$ --- начальная энергия частицы.

Таким образом, пробег определяется скоростью уменьшения энергии
частицы при движении через вещество.


% ============================================================
\section{Многократное рассеяние}
% ============================================================

При прохождении заряженной частицы через вещество она многократно
рассеивается в кулоновских полях атомных ядер.
Каждое отдельное рассеяние обычно происходит на небольшой угол,
однако накопление большого числа таких отклонений приводит
к заметному изменению направления движения частицы.

\begin{rbox}[left=18pt,right=18pt,top=14pt,bottom=10pt]{[]}
    \centering

    \begin{tikzpicture}[
        x=1cm,
        y=1cm,
        >=Latex,
        every node/.style={text=rtext}
    ]

        % ------------------------------------------------------------
        % Оси
        % ------------------------------------------------------------

        \draw[
            ->,
            draw=rtext,
            line width=0.9pt
        ]
        (0,0) -- (7.7,0);

        \draw[
            ->,
            draw=rtext,
            line width=0.9pt
        ]
        (0,0) -- (0,4.8);


        % ------------------------------------------------------------
        % Подписи осей
        % ------------------------------------------------------------

        \node[
            anchor=south,
            font=\small
        ]
        at (-0.05,4.60)
        {$\displaystyle \frac{dN}{dx}$};

        \node[
            anchor=west,
            font=\small
        ]
        at (7.55,-0.03)
        {$\theta_{\min}$};


        % ------------------------------------------------------------
        % Отметка 0.1
        % ------------------------------------------------------------

        \draw[
            draw=rtext,
            line width=0.7pt
        ]
        (-0.10,1.25) -- (0.10,1.25);

        \node[
            anchor=east,
            font=\small
        ]
        at (-0.18,1.25)
        {$0.1$};

        % Пунктир от уровня 0.1 до кривой
        \draw[
            draw=rtext,
            dashed,
            line width=0.65pt
        ]
        (0.10,1.25) -- (2.15,1.25);


        % ------------------------------------------------------------
        % Верхняя отметка 1000
        % ------------------------------------------------------------

        \node[
            anchor=east,
            font=\small
        ]
        at (-0.16,4.15)
        {$1000$};


        % ------------------------------------------------------------
        % Деления по оси theta_min
        % ------------------------------------------------------------

        \draw[
            draw=rtext,
            line width=0.7pt
        ]
        (2.15,-0.10) -- (2.15,0.10);

        \node[
            below=4pt,
            font=\small
        ]
        at (2.15,0)
        {$0.002$};


        \draw[
            draw=rtext,
            line width=0.7pt
        ]
        (6.55,-0.10) -- (6.55,0.10);

        \node[
            below=4pt,
            font=\small
        ]
        at (6.55,0)
        {$0.01$};


        % ------------------------------------------------------------
        % Распределение
        %
        % Резкое падение при малых углах,
        % затем длинный медленно убывающий хвост
        % ------------------------------------------------------------

        \draw[
            draw=rc3,
            line width=1.6pt
        ]
        (0.42,4.35)
        .. controls
            (0.48,3.30)
            and
            (0.67,2.45)
        ..
        (1.05,1.85)
        .. controls
            (1.35,1.47)
            and
            (1.75,1.31)
        ..
        (2.20,1.23)
        .. controls
            (3.10,1.07)
            and
            (4.45,1.00)
        ..
        (6.55,0.97);


        % ------------------------------------------------------------
        % Точка пересечения с уровнем 0.1
        % ------------------------------------------------------------

        \fill[
            rc4
        ]
        (2.15,1.25) circle (1.6pt);

    \end{tikzpicture}

    \captionsetup{
        font=footnotesize,
        labelfont=bf,
        skip=5pt
    }

    \captionof{figure}{
        Качественная зависимость числа рассеяний
        от минимального угла рассеяния $\theta_{\min}$.
        Рассеяния на малые углы происходят значительно чаще,
        а вероятность больших отклонений быстро уменьшается.
    }
    \label{fig:multiple_scattering_theta}
\end{rbox}

Характерная ширина углового распределения определяется выражением

\begin{equation*}
    \sqrt{\langle\theta^2\rangle}
    =
    \frac{13.6}{\beta pc}
    \sqrt{\frac{x'}{X_0}}
    \left[
        1+0.038\ln\left(\frac{x'}{X_0}\right)
    \right].
\end{equation*}

Здесь $p$ --- импульс частицы,
\( \beta=v / c ,  \) 
$x'$ --- пройденная толщина вещества, а $X_0$ ---
\rconcept{радиационная длина} вещества.

При увеличении толщины вещества характерный угол многократного
рассеяния возрастает приблизительно как

\begin{equation*}
    \sqrt{\langle\theta^2\rangle}
    \propto
    \sqrt{\frac{x'}{X_0}}.
\end{equation*}


% ============================================================
\section{Радиационные потери энергии}
% ============================================================

Для электронов при достаточно больших энергиях существенную роль
начинают играть радиационные потери, связанные прежде всего
с тормозным излучением.

Средние радиационные потери энергии можно записать в виде

\begin{equation*}
    \left(
        \frac{dE}{dx}
    \right)_{\mathrm{rad}}
    \approx
    \frac{E}{X_0}.
\end{equation*}

Если учитывать уменьшение самой энергии электрона,

\begin{equation*}
    \frac{dE}{dx}
    =
    -\frac{E}{X_0},
\end{equation*}

то после интегрирования получаем

\begin{equation*}
    E(x)
    =
    E_0
    e^{-x/X_0}.
\end{equation*}

Таким образом, радиационная длина $X_0$ задаёт характерный масштаб,
на котором энергия высокоэнергетического электрона уменьшается
за счёт тормозного излучения.


% ------------------------------------------------------------
\subsection{Тормозное излучение}
% ------------------------------------------------------------

При прохождении электрона в кулоновском поле ядра его траектория
искривляется, вследствие чего электрон испускает фотон.

Соответствующий процесс имеет вид

\begin{equation*}
    e^-+Z
    \rightarrow
    e^-+\gamma+Z.
\end{equation*}

\begin{rbox}[left=18pt,right=18pt,top=14pt,bottom=10pt]{[]}
    \centering

    \begin{tikzpicture}[
        >=Latex,
        every node/.style={text=rtext}
    ]

        % Вершина взаимодействия
        \coordinate (V) at (0,0);

        % ------------------------------------------------------------
        % Входящий электрон
        % ------------------------------------------------------------

        \draw[
            ->,
            draw=rc3,
            line width=1.15pt
        ]
        (-3.2,1.55) -- (V);

        \node[
            above left,
            font=\small
        ]
        at (-2.55,1.25)
        {$e^-$};


        % ------------------------------------------------------------
        % Рассеянный электрон
        % ------------------------------------------------------------

        \draw[
            ->,
            draw=rc3,
            line width=1.15pt
        ]
        (V) -- (2.85,1.25);

        \node[
            above,
            font=\small
        ]
        at (2.35,1.05)
        {$e^-$};


        % ------------------------------------------------------------
        % Реальный тормозной фотон
        % ------------------------------------------------------------

        \draw[
            ->,
            draw=rtext,
            line width=0.9pt,
            decorate,
            decoration={
                snake,
                amplitude=1.7pt,
                segment length=7pt
            }
        ]
        (V) -- (3.55,-0.10);

        \node[
            above,
            font=\small
        ]
        at (3.15,0.05)
        {$\gamma$};


        % ------------------------------------------------------------
        % Виртуальный фотон к ядру
        % ------------------------------------------------------------

        \draw[
            draw=rtext,
            line width=0.85pt,
            decorate,
            decoration={
                snake,
                amplitude=1.5pt,
                segment length=6pt
            }
        ]
        (V) -- (0,-2.15);


        % ------------------------------------------------------------
        % Ядро
        % ------------------------------------------------------------

        \draw[
            draw=rc4,
            line width=1.1pt
        ]
        (0,-2.45) circle (0.30);

        \node[
            font=\small
        ]
        at (0,-2.45)
        {$Z$};

        \node[
            right=7pt,
            font=\small
        ]
        at (0.28,-2.45)
        {ядро};


        % Вершина
        \fill[rc4] (V) circle (1.7pt);

    \end{tikzpicture}

    \captionsetup{
        font=footnotesize,
        labelfont=bf,
        skip=5pt
    }

    \captionof{figure}{
        Тормозное излучение электрона в кулоновском поле ядра.
        При рассеянии на ядре электрон изменяет импульс
        и испускает фотон $\gamma$.
    }
    \label{fig:bremsstrahlung}
\end{rbox}

% ------------------------------------------------------------
\subsection{Критическая энергия}
% ------------------------------------------------------------

При небольших энергиях электрона основную роль играют
ионизационные потери, тогда как при больших энергиях начинают
доминировать радиационные потери.

\rconcept{Критическая энергия} $E_{\text{кр}}$ определяется условием
равенства этих двух механизмов:

\begin{equation*}
    \left(
        \frac{dE}{dx}
    \right)_{\text{ион}}
    =
    \left(
        \frac{dE}{dx}
    \right)_{\text{рад}}.
\end{equation*}

Поскольку

\begin{equation*}
    \left(
        \frac{dE}{dx}
    \right)_{\text{рад}}
    \approx
    \frac{E}{X_0},
\end{equation*}

получаем оценку

\begin{equation*}
    E_{\text{кр}}
    \approx
    \left(
        \frac{dE}{dx}
    \right)_{\text{ион}}
    X_0.
\end{equation*}

Для вещества с атомным номером $Z$ используется приближённое выражение

\begin{equation*}
    E_{\text{кр}}[\mathrm{MeV}]
    \approx
    \frac{610}{Z+1.24}.
\end{equation*}

При

\begin{equation*}
    E\gg E_{\text{кр}}
\end{equation*}

преобладают радиационные потери, тогда как при

\begin{equation*}
    E\ll E_{\text{кр}}
\end{equation*}

основную роль играют ионизационные потери.


% ------------------------------------------------------------
\subsection{Спектр тормозного излучения}
% ------------------------------------------------------------

Энергия электрона при тормозном излучении распределяется
между электроном после взаимодействия и испущенным фотоном.

Число фотонов с энергией $E_{\mathrm{ph}}$ определяется выражением

\begin{equation*}
    \frac{dn}{dE_{\mathrm{ph}}\,dx}
    =
    \frac{1}{X_0E_{\mathrm{ph}}}
    \left(
        \frac{4}{3}
        -\frac{4}{3}y
        +y^2
    \right),
\end{equation*}

где

\begin{equation*}
    y
    =
    \frac{E_{\mathrm{ph}}}{E_e}
\end{equation*}

--- доля энергии электрона, унесённая испущенным фотоном.

Данное выражение справедливо в области полного экранирования,
для которой выполняется условие

\begin{equation*}
    \frac{50\,\mathrm{MeV}}{E_e}
    \frac{y}{1-y}
    Z^{-1/3}
    \ll
    1.
\end{equation*}

При небольших энергиях фотона основная зависимость спектра имеет вид

\begin{equation*}
    \frac{dn}{dE_{\mathrm{ph}}}
    \propto
    \frac{1}{E_{\mathrm{ph}}},
\end{equation*}

поэтому мягких фотонов испускается значительно больше,
чем фотонов с большой энергией.


% ============================================================
\section{Взаимодействие фотонов с веществом}
% ============================================================

При прохождении фотона через вещество возможны различные процессы
взаимодействия. К основным относятся фотоэффект,
комптоновское рассеяние и рождение электрон-позитронных пар.


% ------------------------------------------------------------
\subsection{Рождение электрон-позитронной пары}
% ------------------------------------------------------------

При достаточно большой энергии фотона в кулоновском поле ядра
может происходить рождение электрон-позитронной пары:

\begin{equation*}
    \gamma+Z
    \rightarrow
    e^-+e^++Z.
\end{equation*}

Присутствие ядра необходимо для выполнения закона сохранения импульса.

\begin{rbox}[left=18pt,right=18pt,top=14pt,bottom=10pt]{[]}
    \centering

    \begin{tikzpicture}[
        >=Latex,
        every node/.style={text=rtext},
        %
        photon/.style={
            draw=rtext,
            line width=0.9pt,
            decorate,
            decoration={
                snake,
                amplitude=1.5pt,
                segment length=6pt
            }
        },
        %
        fermion/.style={
            draw=rc3,
            line width=1.15pt
        }
    ]

        % ============================================================
        % ПЕРВАЯ ДИАГРАММА
        %
        % Сначала вдоль фермионной линии находится вершина
        % взаимодействия с виртуальным фотоном ядра,
        % затем вершина поглощения реального фотона.
        % ============================================================

        \begin{scope}[shift={(-3.7,0)}]

            % Вершины
            \coordinate (A) at (-0.75,0.05);
            \coordinate (B) at (0.70,0.48);

            % --------------------------------------------------------
            % Внешний позитрон
            % Стрелка фермионного потока направлена К диаграмме
            % --------------------------------------------------------

            \draw[
                fermion,
                ->
            ]
            (-2.55,-0.28) -- (A);

            \node[
                left,
                font=\small
            ]
            at (-2.52,-0.28)
            {$e^+$};


            % --------------------------------------------------------
            % Внутренняя фермионная линия
            % --------------------------------------------------------

            \draw[
                fermion
            ]
            (A) -- (B);


            % --------------------------------------------------------
            % Выходящий электрон
            % --------------------------------------------------------

            \draw[
                fermion,
                ->
            ]
            (B) -- (2.55,0.17);

            \node[
                right,
                font=\small
            ]
            at (2.55,0.17)
            {$e^-$};


            % --------------------------------------------------------
            % Входящий реальный фотон
            % --------------------------------------------------------

            \draw[
                photon
            ]
            (-0.05,2.02) -- (B);

            \node[
                above left,
                font=\small
            ]
            at (-0.02,1.90)
            {$\gamma$};


            % --------------------------------------------------------
            % Виртуальный фотон от ядра
            % --------------------------------------------------------

            \draw[
                photon
            ]
            (A) -- (-0.75,-1.50);

            \node[
                right,
                font=\scriptsize
            ]
            at (-0.68,-0.80)
            {$\gamma^*$};


            % --------------------------------------------------------
            % Ядро
            % --------------------------------------------------------

            \draw[
                draw=rc4,
                line width=1.0pt
            ]
            (-0.75,-1.80) circle (0.27);

            \node[
                font=\small
            ]
            at (-0.75,-1.80)
            {$Z$};


            % Вершины взаимодействия
            \fill[rc4] (A) circle (1.6pt);
            \fill[rc4] (B) circle (1.6pt);

        \end{scope}


        % ============================================================
        % ВТОРАЯ ДИАГРАММА
        %
        % Теперь порядок вершин обратный:
        % сначала поглощается реальный фотон,
        % затем происходит взаимодействие с полем ядра.
        % ============================================================

        \begin{scope}[shift={(3.7,0)}]

            % Вершины
            \coordinate (C) at (-0.70,0.48);
            \coordinate (D) at (0.75,0.05);

            % --------------------------------------------------------
            % Внешний позитрон
            % --------------------------------------------------------

            \draw[
                fermion,
                ->
            ]
            (-2.55,0.17) -- (C);

            \node[
                left,
                font=\small
            ]
            at (-2.52,0.17)
            {$e^+$};


            % --------------------------------------------------------
            % Внутренняя фермионная линия
            % --------------------------------------------------------

            \draw[
                fermion
            ]
            (C) -- (D);


            % --------------------------------------------------------
            % Выходящий электрон
            % --------------------------------------------------------

            \draw[
                fermion,
                ->
            ]
            (D) -- (2.55,-0.28);

            \node[
                right,
                font=\small
            ]
            at (2.55,-0.28)
            {$e^-$};


            % --------------------------------------------------------
            % Входящий реальный фотон
            % --------------------------------------------------------

            \draw[
                photon
            ]
            (-1.42,2.02) -- (C);

            \node[
                above left,
                font=\small
            ]
            at (-1.38,1.90)
            {$\gamma$};


            % --------------------------------------------------------
            % Виртуальный фотон от ядра
            % --------------------------------------------------------

            \draw[
                photon
            ]
            (D) -- (0.75,-1.50);

            \node[
                right,
                font=\scriptsize
            ]
            at (0.82,-0.80)
            {$\gamma^*$};


            % --------------------------------------------------------
            % Ядро
            % --------------------------------------------------------

            \draw[
                draw=rc4,
                line width=1.0pt
            ]
            (0.75,-1.80) circle (0.27);

            \node[
                font=\small
            ]
            at (0.75,-1.80)
            {$Z$};


            % Вершины взаимодействия
            \fill[rc4] (C) circle (1.6pt);
            \fill[rc4] (D) circle (1.6pt);

        \end{scope}

    \end{tikzpicture}

    \captionsetup{
        font=footnotesize,
        labelfont=bf,
        skip=5pt
    }

    \captionof{figure}{
        Диаграммы рождения электрон-позитронной пары
        фотоном в кулоновском поле ядра.
        Две диаграммы соответствуют двум возможным порядкам
        взаимодействия реального фотона $\gamma$
        и виртуального фотона $\gamma^*$,
        связанного с полем ядра.
    }
    \label{fig:bethe_heitler_pair_production}
\end{rbox}


% ------------------------------------------------------------
\subsection{Фотоэффект}
% ------------------------------------------------------------

При фотоэффекте фотон поглощается связанным электроном,
который в результате покидает атом.

Для сечения фотоэффекта используется зависимость

\begin{equation*}
    \sigma_{\phi}
    =
    \sqrt{\frac{32}{\xi^7}}\,
    \alpha^4 Z^5
    \sigma_{\mathrm{Th}},
\end{equation*}

где $\alpha$ --- постоянная тонкой структуры,
$Z$ --- атомный номер вещества,
а $\sigma_{\mathrm{Th}}$ --- томсоновское сечение.

Безразмерная энергия фотона определяется как

\begin{equation*}
    \xi
    =
    \frac{E_\gamma}{m_e}.
\end{equation*}

Томсоновское сечение равно

\begin{equation*}
    \sigma_{\mathrm{Th}}
    =
    \frac{8}{3}\pi r_e^2
    =
    665\,\mathrm{mb}.
\end{equation*}

Единица площади --- барн:

\begin{equation*}
    1\,\mathrm{b}
    =
    10^{-24}\,\mathrm{cm^2}.
\end{equation*}

Следовательно,

\begin{equation*}
    665\,\mathrm{mb}
    =
    0.665\,\mathrm{b}.
\end{equation*}


% ------------------------------------------------------------
\subsection{Комптоновское рассеяние}
% ------------------------------------------------------------

При комптоновском рассеянии фотон взаимодействует с электроном
и передаёт ему часть своей энергии:

\begin{equation*}
    \gamma+e^-
    \rightarrow
    \gamma+e^-.
\end{equation*}

\begin{rbox}[left=18pt,right=18pt,top=14pt,bottom=10pt]{[]}
    \centering

    \begin{tikzpicture}[
        >=Latex,
        every node/.style={text=rtext},
        %
        fermion/.style={
            draw=rc3,
            line width=1.15pt
        },
        %
        photon/.style={
            draw=rtext,
            line width=0.9pt,
            decorate,
            decoration={
                snake,
                amplitude=1.5pt,
                segment length=6pt
            }
        }
    ]

        % ============================================================
        % Прямая диаграмма
        % ============================================================

        \begin{scope}[shift={(-3.8,0)}]

            % Вершины
            \coordinate (A) at (-0.9,0);
            \coordinate (B) at (0.9,0);

            % --------------------------------------------------------
            % Электронная линия
            % --------------------------------------------------------

            \draw[
                fermion,
                ->
            ]
            (-2.7,0) -- (A);

            \draw[
                fermion,
                ->
            ]
            (A) -- (B);

            \draw[
                fermion,
                ->
            ]
            (B) -- (2.7,0);

            % Подписи электронов
            \node[
                below,
                font=\small
            ]
            at (-2.25,0)
            {$e^-$};

            \node[
                below,
                font=\small
            ]
            at (2.25,0)
            {$e^-$};


            % --------------------------------------------------------
            % Входящий фотон
            % --------------------------------------------------------

            \draw[
                photon,
                ->
            ]
            (-2.15,1.85) -- (A);

            \node[
                above left,
                font=\small
            ]
            at (-1.85,1.55)
            {$E_\gamma$};


            % --------------------------------------------------------
            % Выходящий фотон
            % --------------------------------------------------------

            \draw[
                photon,
                ->
            ]
            (B) -- (2.20,1.85);

            \node[
                above right,
                font=\small
            ]
            at (1.85,1.52)
            {$E'_\gamma$};


            % --------------------------------------------------------
            % Вершины
            % --------------------------------------------------------

            \fill[rc4] (A) circle (1.6pt);
            \fill[rc4] (B) circle (1.6pt);


            % Подпись диаграммы
            \node[
                below=13pt,
                font=\small
            ]
            at (0,-0.15)
            {прямая};

        \end{scope}


        % ============================================================
        % Перекрёстная диаграмма
        % ============================================================

        \begin{scope}[shift={(3.8,0)}]

            % Вершины
            \coordinate (C) at (-0.9,0);
            \coordinate (D) at (0.9,0);

            % --------------------------------------------------------
            % Электронная линия
            % --------------------------------------------------------

            \draw[
                fermion,
                ->
            ]
            (-2.7,0) -- (C);

            \draw[
                fermion,
                ->
            ]
            (C) -- (D);

            \draw[
                fermion,
                ->
            ]
            (D) -- (2.7,0);

            % Подписи электронов
            \node[
                below,
                font=\small
            ]
            at (-2.25,0)
            {$e^-$};

            \node[
                below,
                font=\small
            ]
            at (2.25,0)
            {$e^-$};


            % --------------------------------------------------------
            % Входящий фотон
            %
            % Идёт к ПРАВОЙ вершине
            % --------------------------------------------------------

            \draw[
                photon,
                ->
            ]
            (-2.15,1.85)
            -- (D);

            \node[
                above left,
                font=\small
            ]
            at (-1.85,1.55)
            {$E_\gamma$};


            % --------------------------------------------------------
            % Выходящий фотон
            %
            % Выходит из ЛЕВОЙ вершины
            % --------------------------------------------------------

            \draw[
                photon,
                ->
            ]
            (C)
            -- (2.20,1.85);

            \node[
                above right,
                font=\small
            ]
            at (1.85,1.52)
            {$E'_\gamma$};


            % --------------------------------------------------------
            % Вершины
            % --------------------------------------------------------

            \fill[rc4] (C) circle (1.6pt);
            \fill[rc4] (D) circle (1.6pt);


            % Подпись диаграммы
            \node[
                below=13pt,
                font=\small
            ]
            at (0,-0.15)
            {перекрёстная};

        \end{scope}

    \end{tikzpicture}

    \captionsetup{
        font=footnotesize,
        labelfont=bf,
        skip=5pt
    }

    \captionof{figure}{
        Диаграммы Комптон-рассеяния
        $e^-+\gamma\rightarrow e^-+\gamma$.
        Слева показана прямая диаграмма,
        справа --- перекрёстная.
    }
    \label{fig:compton_diagrams}
\end{rbox}

Введём безразмерную энергию фотона

\begin{equation*}
    \xi
    =
    \frac{E_\gamma}{m_e}.
\end{equation*}

Если после взаимодействия фотон рассеивается на угол $\theta_\gamma$,
то отношение его конечной и начальной энергий равно

\begin{equation*}
    \frac{E_\gamma'}{E_\gamma}
    =
    \frac{1}{
        1+\xi(1-\cos\theta_\gamma)
    }.
\end{equation*}

Следовательно,

\begin{equation*}
    E_\gamma'
    =
    \frac{E_\gamma}{
        1+\xi(1-\cos\theta_\gamma)
    }.
\end{equation*}

При увеличении угла рассеяния энергия конечного фотона уменьшается,
а всё большая доля начальной энергии передаётся электрону.


% ------------------------------------------------------------
\subsection{Сечение рождения пары}
% ------------------------------------------------------------

При высоких энергиях одним из основных механизмов поглощения фотонов
становится рождение электрон-позитронных пар в поле ядра.

\begin{rbox}[left=18pt,right=18pt,top=14pt,bottom=10pt]{[]}
    \centering

    \begin{tikzpicture}[
        >=Latex,
        every node/.style={text=rtext},
        photon/.style={
            draw=rtext,
            line width=0.9pt,
            decorate,
            decoration={
                snake,
                amplitude=1.5pt,
                segment length=6pt
            }
        },
        charged/.style={
            draw=rc3,
            line width=1.15pt,
            -> 
        }
    ]

        % Вершина рождения пары
        \coordinate (V) at (0,0);

        % Входящий фотон
        \draw[
            photon,
            ->
        ]
        (-3.1,0) -- (V);

        \node[
            above,
            font=\small
        ]
        at (-2.45,0.15)
        {$\gamma$};

        % Электрон
        \draw[
            charged
        ]
        (V) -- (2.75,1.25);

        \node[
            above right,
            font=\small
        ]
        at (2.35,1.05)
        {$e^-$};

        % Позитрон
        \draw[
            charged
        ]
        (V) -- (2.75,-1.25);

        \node[
            below right,
            font=\small
        ]
        at (2.35,-1.05)
        {$e^+$};

        % Виртуальный фотон от ядра
        \draw[
            photon
        ]
        (V) -- (0,-2.0);

        \node[
            right,
            font=\scriptsize
        ]
        at (0.10,-1.15)
        {$\gamma^*$};

        % Ядро
        \draw[
            draw=rc4,
            line width=1.05pt
        ]
        (0,-2.35) circle (0.30);

        \node[
            font=\small
        ]
        at (0,-2.35)
        {$Z$};

        \node[
            right=7pt,
            font=\small
        ]
        at (0.30,-2.35)
        {ядро};

        % Вершина
        \fill[rc4]
        (V) circle (1.7pt);

    \end{tikzpicture}

    \captionsetup{
        font=footnotesize,
        labelfont=bf,
        skip=5pt
    }

    \captionof{figure}{
        Схематическое изображение рождения
        электрон-позитронной пары фотоном
        в кулоновском поле ядра.
    }
    \label{fig:pair_production_scheme}
\end{rbox}

Сечение рождения пары записывается в виде

\begin{equation*}
    \sigma_n^{\text{пар}}
    =
    4\alpha r_e^2
    \left(
        \frac{7}{9}
        \lg\frac{183}{Z^{1/3}}
        -
        \frac{1}{54}
    \right).
\end{equation*}

Вероятность рождения пары при прохождении малого участка вещества
толщины $dx$ составляет

\begin{equation*}
    dW
    \approx
    \frac{7}{9}
    \frac{dx}{X_0}.
\end{equation*}

Таким образом, характерная длина взаимодействия фотона
при рождении пар также определяется радиационной длиной $X_0$.


% ============================================================
\section{Электрон-фотонный ливень}
% ============================================================

При достаточно большой энергии первоначальной частицы
тормозное излучение и рождение электрон-позитронных пар
могут многократно чередоваться.

Электрон или позитрон испускает тормозной фотон,

\begin{equation*}
    e^\pm
    \rightarrow
    e^\pm+\gamma,
\end{equation*}

после чего фотон рождает новую электрон-позитронную пару:

\begin{equation*}
    \gamma
    \rightarrow
    e^-+e^+.
\end{equation*}

Новые электроны и позитроны снова испускают фотоны,
которые, в свою очередь, рождают новые пары.
Возникает \rconcept{электрон-фотонный ливень}.

\begin{rbox}[left=18pt,right=18pt,top=14pt,bottom=10pt]{[]}
    \centering

    \begin{tikzpicture}[
        >=Latex,
        every node/.style={text=rtext},
        photon/.style={
            draw=rtext,
            line width=0.85pt,
            decorate,
            decoration={
                snake,
                amplitude=1.35pt,
                segment length=5.5pt
            }
        },
        charged/.style={
            draw=rc3,
            line width=1.05pt,
            ->
        }
    ]

        % ============================================================
        % Первый фотон
        % ============================================================

        \coordinate (A) at (-3.3,0);
        \coordinate (B) at (-1.9,0);

        \draw[
            photon,
            ->
        ]
        (A) -- (B);

        \node[
            above,
            font=\small
        ]
        at (-2.65,0.15)
        {$\gamma$};


        % ============================================================
        % Первая пара
        % ============================================================

        \coordinate (C1) at (-0.45,1.10);
        \coordinate (C2) at (-0.45,-1.10);

        \draw[charged]
        (B) -- (C1);

        \draw[charged]
        (B) -- (C2);

        \node[
            above,
            font=\scriptsize
        ]
        at (-0.82,0.77)
        {$e^-$};

        \node[
            below,
            font=\scriptsize
        ]
        at (-0.82,-0.77)
        {$e^+$};


        % ============================================================
        % Тормозные фотоны
        % ============================================================

        \coordinate (D1) at (0.75,1.65);
        \coordinate (D2) at (0.75,0.50);
        \coordinate (D3) at (0.75,-0.50);
        \coordinate (D4) at (0.75,-1.65);

        % Продолжение заряженных частиц
        \draw[charged]
        (C1) -- (1.15,1.30);

        \draw[charged]
        (C2) -- (1.15,-1.30);

        % Фотоны от e-
        \draw[
            photon,
            ->
        ]
        (C1) -- (D1);

        \draw[
            photon,
            ->
        ]
        (C1) -- (D2);

        % Фотоны от e+
        \draw[
            photon,
            ->
        ]
        (C2) -- (D3);

        \draw[
            photon,
            ->
        ]
        (C2) -- (D4);


        % ============================================================
        % Второе поколение пар
        % ============================================================

        % Верхний фотон
        \coordinate (E11) at (2.20,2.10);
        \coordinate (E12) at (2.20,1.35);

        \draw[charged]
        (D1) -- (E11);

        \draw[charged]
        (D1) -- (E12);

        % Второй сверху
        \coordinate (E21) at (2.20,0.85);
        \coordinate (E22) at (2.20,0.15);

        \draw[charged]
        (D2) -- (E21);

        \draw[charged]
        (D2) -- (E22);

        % Второй снизу
        \coordinate (E31) at (2.20,-0.15);
        \coordinate (E32) at (2.20,-0.85);

        \draw[charged]
        (D3) -- (E31);

        \draw[charged]
        (D3) -- (E32);

        % Нижний
        \coordinate (E41) at (2.20,-1.35);
        \coordinate (E42) at (2.20,-2.10);

        \draw[charged]
        (D4) -- (E41);

        \draw[charged]
        (D4) -- (E42);


        % ============================================================
        % Ещё несколько тормозных фотонов
        % ============================================================

        \draw[
            photon,
            ->
        ]
        (E11) -- (3.20,2.45);

        \draw[
            photon,
            ->
        ]
        (E12) -- (3.20,1.55);

        \draw[
            photon,
            ->
        ]
        (E21) -- (3.20,1.05);

        \draw[
            photon,
            ->
        ]
        (E22) -- (3.20,0.25);

        \draw[
            photon,
            ->
        ]
        (E31) -- (3.20,-0.25);

        \draw[
            photon,
            ->
        ]
        (E32) -- (3.20,-1.05);

        \draw[
            photon,
            ->
        ]
        (E41) -- (3.20,-1.55);

        \draw[
            photon,
            ->
        ]
        (E42) -- (3.20,-2.45);


        % ============================================================
        % Условные границы поколений
        % ============================================================

        \draw[
            draw=rtext,
            opacity=0.25,
            dashed,
            line width=0.6pt
        ]
        (-1.90,-2.65) -- (-1.90,2.65);

        \draw[
            draw=rtext,
            opacity=0.25,
            dashed,
            line width=0.6pt
        ]
        (0.75,-2.65) -- (0.75,2.65);

        \draw[
            draw=rtext,
            opacity=0.25,
            dashed,
            line width=0.6pt
        ]
        (2.20,-2.65) -- (2.20,2.65);


        % Подписи характерного масштаба
        \node[
            below,
            font=\scriptsize
        ]
        at (-0.58,-2.65)
        {$X_0$};

        \node[
            below,
            font=\scriptsize
        ]
        at (1.48,-2.65)
        {$2X_0$};

        \node[
            below,
            font=\scriptsize
        ]
        at (2.85,-2.65)
        {$3X_0$};

    \end{tikzpicture}

    \captionsetup{
        font=footnotesize,
        labelfont=bf,
        skip=5pt
    }

    \captionof{figure}{
        Схематическое развитие электрон-фотонного ливня.
        Рождение электрон-позитронных пар и тормозное излучение
        последовательно увеличивают число вторичных частиц.
    }
    \label{fig:electromagnetic_shower}
\end{rbox}


% ------------------------------------------------------------
\subsection{Развитие электрон-фотонного ливня}
% ------------------------------------------------------------

В простейшей каскадной модели после прохождения каждого характерного
участка порядка радиационной длины число частиц приблизительно
удваивается.

Поэтому число частиц на глубине $x$ можно оценить как

\begin{equation*}
    N_s
    \sim
    2^{x/X_0}.
\end{equation*}

По мере развития ливня энергия первоначальной частицы распределяется
между всё большим числом вторичных частиц.

Средняя энергия одной частицы составляет

\begin{equation*}
    \overline{E}_s
    =
    \frac{E_0}{N_s},
\end{equation*}

где $E_0$ --- начальная энергия.

Каскадное размножение продолжается до тех пор, пока энергия
электронов и позитронов достаточно велика и радиационные процессы
преобладают над ионизационными потерями.

Когда средняя энергия частиц становится порядка критической,

\begin{equation*}
    \overline{E}_s
    \sim
    E_{\text{кр}},
\end{equation*}

дальнейшее размножение ливня прекращается.

В максимуме ливня можно положить

\begin{equation*}
    \frac{E_0}{N_s}
    =
    E_{\text{кр}},
\end{equation*}

откуда

\begin{equation*}
    N_s
    =
    \frac{E_0}{E_{\text{кр}}}.
\end{equation*}

С другой стороны,

\begin{equation*}
    N_s
    \sim
    2^{x/X_0}.
\end{equation*}

Поэтому в точке максимального развития ливня

\begin{equation*}
    2^{x_{\max}/X_0}
    =
    \frac{E_0}{E_{\text{кр}}}.
\end{equation*}

Логарифмируя это выражение по основанию $2$, получаем

\begin{equation*}
    \frac{x_{\max}}{X_0}
    =
    \log_2
    \frac{E_0}{E_{\text{кр}}}.
\end{equation*}

Следовательно,

\begin{equation*}
    x_{\max}
    =
    X_0
    \log_2
    \frac{E_0}{E_{\text{кр}}}.
\end{equation*}

Таким образом, глубина максимума электрон-фотонного ливня
увеличивается логарифмически с ростом начальной энергии частицы.

\begin{rbox}[left=18pt,right=18pt,top=14pt,bottom=10pt]{[]}
    \centering

    \begin{tikzpicture}[
        x=1cm,
        y=1cm,
        >=Latex,
        every node/.style={text=rtext}
    ]

        % ============================================================
        % Оси
        % ============================================================

        \draw[
            ->,
            draw=rtext,
            line width=0.9pt
        ]
        (0,0) -- (7.8,0);

        \draw[
            ->,
            draw=rtext,
            line width=0.9pt
        ]
        (0,0) -- (0,4.7);


        % Подписи осей
        \node[
            anchor=west,
            font=\small
        ]
        at (7.55,-0.03)
        {$x$};

        \node[
            anchor=south,
            font=\small
        ]
        at (-0.05,4.50)
        {$N(x)$};


        % ============================================================
        % Продольный профиль
        % ============================================================

        \draw[
            draw=rc3,
            line width=1.65pt
        ]
        (0.20,0.12)
        .. controls
            (0.75,0.15)
            and
            (1.25,0.55)
        ..
        (1.85,1.55)
        .. controls
            (2.45,2.70)
            and
            (3.00,3.75)
        ..
        (3.75,3.85)
        .. controls
            (4.40,3.85)
            and
            (4.70,3.15)
        ..
        (4.95,2.20)
        .. controls
            (5.20,1.20)
            and
            (5.65,0.45)
        ..
        (6.45,0.18);


        % ============================================================
        % Максимум
        % ============================================================

        \coordinate (M) at (3.75,3.85);

        \draw[
            draw=rtext,
            dashed,
            opacity=0.65,
            line width=0.7pt
        ]
        (M) -- (3.75,0);

        \fill[
            rc4
        ]
        (M) circle (1.8pt);


        % Метка x_max
        \draw[
            draw=rtext,
            line width=0.7pt
        ]
        (3.75,-0.10) -- (3.75,0.10);

        \node[
            below=4pt,
            font=\small
        ]
        at (3.75,0)
        {$x_{\max}$};


        % Подпись максимума
        \node[
            above,
            font=\small
        ]
        at (M)
        {$N_{\max}$};

    \end{tikzpicture}

    \captionsetup{
        font=footnotesize,
        labelfont=bf,
        skip=5pt
    }

    \captionof{figure}{
        Продольное развитие электрон-фотонного ливня.
        Число частиц возрастает до максимума на глубине
        $x_{\max}$, после чего уменьшается.
    }
    \label{fig:shower_longitudinal_profile}
\end{rbox}

\end{document}
